\documentclass[11pt]{article}
\usepackage[margin=1in]{geometry}
\usepackage{amsmath,amsthm,amssymb,mathtools}
\usepackage{enumitem,hyperref,mdframed}

\newtheorem{theorem}{Theorem}[section]
\newtheorem{lemma}[theorem]{Lemma}
\newtheorem{proposition}[theorem]{Proposition}
\newtheorem{corollary}[theorem]{Corollary}
\theoremstyle{definition}
\newtheorem{definition}[theorem]{Definition}
\newtheorem{remark}[theorem]{Remark}

\newcommand{\eps}{\varepsilon}
\newcommand{\R}{\mathbb{R}}
\newcommand{\dd}{\mathrm{d}}
\newcommand{\ddt}{\frac{\mathrm{d}}{\mathrm{d}t}}
\newcommand{\ur}{u^{r}}
\newcommand{\uz}{u^{z}}
\newcommand{\ut}{u^{\theta}}
\newcommand{\Ph}{\Phi}
\newcommand{\Ga}{\Gamma}
\newcommand{\Lfour}{\mathcal{L}_{4}}
\newcommand{\gradfive}{\nabla_{5}}
\newcommand{\divfive}{\mathrm{div}_{5}}
\newcommand{\Hdot}{\dot{H}}
\newcommand{\Lint}{L^{\infty}}
\newcommand{\norm}[2]{\|#1\|_{#2}}
\newcommand{\abs}[1]{\lvert#1\rvert}

\title{Phi-renormalization for axisymmetric Navier--Stokes with swirl:\\
the five-dimensional energy identity and the critical strain barrier}
\author{Jonathan Simons\\
Prime Field Technologies LLC, Savannah, Georgia}
\date{22 August 2026}

\begin{document}
\maketitle

\begin{abstract}
This paper continues the author's May 2026 Phi-renormalization preprint
(Zenodo concept DOI
\href{https://doi.org/10.5281/zenodo.20405404}{10.5281/zenodo.20405404},
file DOI
\href{https://doi.org/10.5281/zenodo.20405405}{10.5281/zenodo.20405405}).
Kept from that deposit: the substitutions \(\Ga=r\ut\), \(\Ph=\ut/r\),
and the algebraic identity
\(\frac1{r^4}\partial_z(\Ga^2)=\partial_z(\Ph^2)\).
Withdrawn as load-bearing for the classical system: the line that
bounds \(\partial_z(\Ph^2)\) in \(L^2\) by placing \(\Ph\) in \(\Lint\)
``by Sobolev''; the Gronwall-free convergence sketch; and every
prime-lattice / \(Q_6\) / spectral-clock statement.

The \(\Ph\)-energy is taken in the measure \(r^3\,\dd r\,\dd z\),
in which \(\Lfour=\partial_{rr}+(3/r)\partial_r+\partial_{zz}\) is
self-adjoint and the axis boundary term of the \(r\,\dd r\,\dd z\)
calculus is absent. With radial and axial cutoffs displayed and then
removed, the headline identity is
\begin{equation*}
\tfrac12\ddt\norm{\Ph}{L^2(r^3)}^2
+\nu\norm{\gradfive\Ph}{L^2(r^3)}^2
+\eps\norm{\Ph}{\Hdot^{1.3}(r^3)}^2
=-\int\frac{\ur}{r}\,\Ph^2\,r^3\,\dd r\,\dd z.
\tag{\ensuremath{*}}
\end{equation*}
Gronwall closes on \((*)\) if and only if
\(\displaystyle\sup_{\eps\in(0,1]}\int_0^T\norm{\ur_\eps/r}{\Lint}\,\dd t<\infty\).
That bound---called (B1) below---is open. It is equivalent in difficulty
to global regularity for the axisymmetric-with-swirl class. It is the
theorem, not a remaining \(\delta\)-loss.
Two pillars are proved: the maximum principle for \(\Ga\), and global
smoothness of the \(\eps\)-system.
No classical large-data regularity is claimed.
\end{abstract}

\noindent\textit{MSC 2020.} 35Q30, 76D03, 35B65.\\
\textit{Keywords.} Navier--Stokes, axisymmetric swirl, circulation,
five-dimensional lift, critical strain.

\tableofcontents
\vspace{0.75em}

\begin{mdframed}[linewidth=0.9pt]
\textbf{Status for the referee.}
Proved here: the \(\Ph\)-identity; the equation for \(\Ga\) with every
term written; the maximum principle for \(\Ga\); the equation for \(\Ph\)
with the identity \((\Delta-r^{-2})(r\Ph)=r\Lfour\Ph\) expanded;
the localized and global forms of \eqref{eq:star} in \(r^3\,\dd r\,\dd z\),
including \(\divfive(u)=2\ur/r\) and the cutoff remainders;
criticality of the strain pairing by an explicit NS scaling count;
global smoothness and the energy inequality for the \(\eps\)-system.
Open: the uniform strain bound \eqref{eq:B1}.
Not claimed: classical unaugmented global regularity;
Gronwall-free \(\eps\to 0\) on a possibly singular interval;
any prime-lattice, \(Q_6\), Triple Lock, spectral-clock, or Clay statement.
\end{mdframed}

\section{Introduction}

Axisymmetric Navier--Stokes without swirl is globally regular
\cite{Lady1968,UY1968}.
With swirl the large-data problem is open.
The centrifugal source \(\frac1{r^4}\partial_z(\Ga^2)\) in the
meridional vorticity equation is the classical obstruction
\cite{Danchin2007,CSTY2008,CSTY2009,LeiZhang2011}.

The substitution \(\Ph=\ut/r=\Ga/r^2\) rewrites that source as
\(\partial_z(\Ph^2)\). This is algebra, not a bound
(Theorem~\ref{thm:identity}).
The analytic difficulty moves onto \(\Ph\) and onto the radial strain
\(\ur/r\). The identity does not place \(\Ph\) in \(\Lint\), and it
does not close a Beale--Kato--Majda integral.

The operator \(\partial_{rr}+(3/r)\partial_r+\partial_{zz}\) is the radial
Laplacian of \(\R^4\times\R\), so the natural measure for \(\Ph\) is
\(r^3\,\dd r\,\dd z\). This is a change of variables in the
three-dimensional problem, not a change of dimension: no physics leaves
\(\R^3\), and every statement can be rewritten in \(r\,\dd r\,\dd z\)
at the cost of an explicit axis boundary term. The higher-dimensional
language is bookkeeping.

Regularity theorems in the swirl class assume a Type~I / strain bound:
\(\abs{\ur}\le C/r\), or
\(\int_0^T\norm{\ur/r}{\Lint}\,\dd t<\infty\),
or comparable control on \(r\ut\)
\cite{CSTY2008,CSTY2009,ChenFangZhang2017}.
The assumption has not been removed.
This paper derives the energy identity that makes that assumption the
only remaining hypothesis on the \(\Ph\)-side, and records why
interpolation cannot buy it.

Announcements of 2026 swirl and full-system closures include
\cite{ShahApr,ShahMay,ShahJun}.
They are cited and not used. In particular the \((A,W)\)
circulation-gradient construction and the axis-Hardy formula of
\cite{ShahJun} are not imported and are not claimed as the author's.
A short, checkable identity plus an accurately named obstruction is
the document that can survive a referee. An unchecked closure cannot.

\section{Public record}
\label{sec:record}

The May 2026 files are public:
\begin{itemize}[leftmargin=1.2em]
\item concept DOI
\href{https://doi.org/10.5281/zenodo.20405404}{10.5281/zenodo.20405404};
\item latest file DOI
\href{https://doi.org/10.5281/zenodo.20405405}{10.5281/zenodo.20405405}
(27 May 2026, \texttt{PhiRenorm\_TrackB.pdf});
\item short companion
\href{https://doi.org/10.5281/zenodo.20405597}{10.5281/zenodo.20405597}
(\texttt{Simons\_PhiRenorm\_Axisymmetric.tex}).
\end{itemize}
A later ``New version'' on the concept record preserves that timestamp.
This note does not replace it. It restricts what that timestamp may be
asked to carry.

What stands, from the public text:
\begin{itemize}[leftmargin=1.2em]
\item \(\Ga=r\ut\) and \(\Ph=\Ga/r^2=\ut/r\) (classical notation);
\item \(\frac1{r^4}\partial_z(\Ga^2)=\partial_z(\Ph^2)\) on \(\{r>0\}\)
(algebra);
\item the same factorization for the difference
\(\Ga_\eps^2-\Ga^2\).
\end{itemize}

What does not stand, as a proof of classical regularity:
\begin{itemize}[leftmargin=1.2em]
\item ``\(\norm{\partial_z(\Ph^2)}{L^2}\le 2\norm{\Ph}{\infty}\norm{\partial_z\Ph}{L^2}\)
by Sobolev on a bounded domain.''
This assumes the \(\Lint\) control that continuation is trying to prove.
\item The Gronwall-free rate \(u^\eps\to u\) at
\(O(\eps^{4/(\beta+2)})\) via a multi-component stability energy.
The cancellation that is supposed to kill the Gronwall driver is a
sketch, not an identity.
\item Global \(C^\infty\) for a \(Q_1\)-augmented system, sold as if it
settled the classical problem. Extra parabolicity is a different PDE.
\item Prime-lattice, \(Q_6\), spectral-clock, and Clay packaging.
Those statements are not used and are not restated as theorems.
\end{itemize}

The May open-problem box already marked classical unaugmented swirl as
open. The present paper keeps that honesty and names the remaining
bound.

\section{Notation and function spaces}

Cylindrical coordinates \((r,\theta,z)\) on \(\R^3\),
\(r=\sqrt{x_1^2+x_2^2}\). An axisymmetric velocity is
\[
  u=\ur(r,z,t)\,e_r+\ut(r,z,t)\,e_\theta+\uz(r,z,t)\,e_z.
\]
Incompressibility is
\begin{equation}
  \label{eq:div3}
  \partial_r(r\ur)+\partial_z(r\uz)=0.
\end{equation}
Smoothness on the axis forces \(\ut=\ur=0\) at \(r=0\) and
\(\ut=O(r)\), \(\ur=O(r)\) as \(r\to 0\). In particular
\(\Ga=O(r^2)\) and \(\Ph\) extends continuously to \(r=0\) by
\(\Ph(0,z)=\partial_r\ut(0,z)\).

\begin{definition}[Circulation and intensive swirl]
\[
  \Ga:=r\ut,\qquad \Ph:=\frac{\ut}{r}=\frac{\Ga}{r^2}.
\]
\end{definition}

Two \(L^2\) spaces are used and are not interchangeable.
\begin{itemize}[leftmargin=1.2em]
\item \(\Ga\) lives in \(L^2(r\,\dd r\,\dd z)\), the physical
three-dimensional energy measure in the meridional half-plane.
\item \(\Ph\) lives in \(L^2(r^3\,\dd r\,\dd z)\):
\[
  \norm{\Ph}{L^2(r^3)}^2
  :=\int_0^\infty\int_{\R}\Ph^2\,r^3\,\dd z\,\dd r.
\]
\end{itemize}
Write \(\gradfive\Ph:=(\partial_r\Ph,\partial_z\Ph)\) and
\[
  \divfive(u)
  :=r^{-3}\bigl(\partial_r(r^3\ur)+\partial_z(r^3\uz)\bigr).
\]
The operator
\[
  \Lfour:=\partial_{rr}+\frac3r\partial_r+\partial_{zz}
\]
satisfies
\begin{equation}
  \label{eq:L4div}
  \partial_{rr}+\frac3r\partial_r
  =r^{-3}\partial_r\bigl(r^3\partial_r\,\cdot\,\bigr).
\end{equation}
Thus \(\Lfour\) is symmetric on \(L^2(r^3\,\dd r\,\dd z)\) for functions
with \(\Ph,\partial_r\Ph\) bounded at \(r=0\) and decaying at infinity.

The \(\eps\)-system uses fractional dissipation of order \(1.3\).
Testing \(\eps(-\Delta)^{1.3}u\) against \(u\) produces
\(\eps\norm{u}{\Hdot^{1.3}}^2\). The number \(2.6\) is the composition
order of the operator, not the energy index. Duality is
\begin{equation}
  \label{eq:dual}
  \bigl\langle(-\Delta)^{1.3}u,\varphi\bigr\rangle
  =\bigl\langle(-\Delta)^{0.65}u,(-\Delta)^{0.65}\varphi\bigr\rangle,
\end{equation}
hence
\(\abs{\langle\eps(-\Delta)^{1.3}u,\varphi\rangle}
\le\eps\norm{u}{\Hdot^{1.3}}\norm{\varphi}{\Hdot^{1.3}}\).

\section{The algebraic identity}

\begin{theorem}[Phi-renormalization]
\label{thm:identity}
On \(\{r>0\}\),
\begin{equation}
  \label{eq:id}
  \frac1{r^4}\partial_z(\Ga^2)=\partial_z(\Ph^2).
\end{equation}
If \(\Ga_\eps=r^2\Ph_\eps\), then
\(\frac1{r^4}\partial_z(\Ga_\eps^2-\Ga^2)
=\partial_z\bigl((\Ph_\eps+\Ph)(\Ph_\eps-\Ph)\bigr)\).
\end{theorem}

\begin{proof}
\(\Ga=r^2\Ph\), so \(\Ga^2=r^4\Ph^2\). Since \(\partial_z(r^4)=0\),
\[
  \frac1{r^4}\partial_z(\Ga^2)
  =\frac1{r^4}\partial_z(r^4\Ph^2)
  =\partial_z(\Ph^2).
\]
The difference identity is the same factorization:
\(\Ga_\eps^2-\Ga^2=r^4(\Ph_\eps+\Ph)(\Ph_\eps-\Ph)\).
\end{proof}

\begin{remark}[Withdrawn May line]
\label{rem:withdrawn}
The identity does not bound \(\Ph\) in \(\Lint\).
The May 2026 estimate
\(\norm{\partial_z(\Ph^2)}{L^2}\le 2\norm{\Ph}{\infty}\norm{\partial_z\Ph}{L^2}\)
``by Sobolev'' assumes the control classical continuation seeks.
It is not used below. Hardy is bypassed for the \emph{rewrite} of the
source; the strain pairing returns the difficulty.
\end{remark}

\section{Pillar I: the maximum principle for \(\Gamma\)}

The azimuthal Navier--Stokes component is
\begin{equation}
  \label{eq:utheta}
  \partial_t\ut+\ur\partial_r\ut+\uz\partial_z\ut+\frac{\ur}{r}\ut
  =\nu\Bigl(\Delta-\frac1{r^2}\Bigr)\ut,
\end{equation}
with \(\Delta=\partial_{rr}+r^{-1}\partial_r+\partial_{zz}\).

\begin{lemma}[Circulation equation]
\label{lem:Gamma-eq}
If \(u\) is a smooth axisymmetric solution, then
\begin{equation}
  \label{eq:Gamma}
  \partial_t\Ga+\ur\partial_r\Ga+\uz\partial_z\Ga
  =\nu\Bigl(\partial_{rr}-\frac1r\partial_r+\partial_{zz}\Bigr)\Ga.
\end{equation}
There is no zeroth-order stretching of the wrong sign.
\end{lemma}

\begin{proof}
\(\partial_t\Ga=r\partial_t\ut\),
\(\partial_r\Ga=\ut+r\partial_r\ut\),
\(\partial_z\Ga=r\partial_z\ut\).
Multiply \eqref{eq:utheta} by \(r\):
\[
  r\partial_t\ut+r\ur\partial_r\ut+r\uz\partial_z\ut+\ur\ut
  =\nu r\Bigl(\Delta-r^{-2}\Bigr)\ut.
\]
The left side is
\[
  \partial_t\Ga+\ur(\partial_r\Ga-\ut)+\uz\partial_z\Ga+\ur\ut
  =\partial_t\Ga+\ur\partial_r\Ga+\uz\partial_z\Ga.
\]
For the viscous side, write \(\ut=\Ga/r\). Then
\begin{align*}
\partial_r\ut&=r^{-1}\partial_r\Ga-r^{-2}\Ga,\\
\partial_{rr}\ut&=r^{-1}\partial_{rr}\Ga-2r^{-2}\partial_r\Ga+2r^{-3}\Ga,\\
\partial_{zz}\ut&=r^{-1}\partial_{zz}\Ga.
\end{align*}
Hence
\begin{align*}
\bigl(\Delta-r^{-2}\bigr)\ut
&=r^{-1}\partial_{rr}\Ga-2r^{-2}\partial_r\Ga+2r^{-3}\Ga
 +r^{-1}\bigl(r^{-1}\partial_r\Ga-r^{-2}\Ga\bigr)
 +r^{-1}\partial_{zz}\Ga-r^{-3}\Ga\\
&=r^{-1}\partial_{rr}\Ga-r^{-2}\partial_r\Ga+r^{-1}\partial_{zz}\Ga.
\end{align*}
Multiplication by \(r\) gives the right-hand side of \eqref{eq:Gamma}.
\end{proof}

\begin{theorem}[Circulation maximum principle]
\label{thm:gamma}
Let \(u\) be a smooth axisymmetric solution of Navier--Stokes, or of
the \(\eps\)-system \eqref{eq:eps}, on \([0,T]\). Then
\[
  \norm{\Ga(t)}{\Lint}\le\norm{\Ga_0}{\Lint}
  \qquad\text{for all }t\in[0,T].
\]
No smallness is required.
\end{theorem}

\begin{proof}
Equation \eqref{eq:Gamma} is a linear scalar advection--diffusion equation
for \(\Ga\). At an interior spatial maximum of \(\Ga(\,\cdot\,,t)\) one
has \(\nabla\Ga=0\) and
\(\partial_{rr}\Ga-r^{-1}\partial_r\Ga+\partial_{zz}\Ga\le 0\),
hence \(\partial_t\Ga\le 0\). The same holds for \(-\Ga\).
On the axis, smoothness forces \(\Ga=O(r^2)\), so \(\Ga|_{r=0}=0\):
an interior maximum of \(\abs{\Ga}\) cannot hide at \(r=0\) unless
\(\Ga\equiv 0\).
The \(\eps\)-system adds a nonpositive fractional term at a maximum
(the Fourier symbol of \((-\Delta)^{1.3}\) is positive).
Thus \(\norm{\Ga(t)}{\Lint}\) is nonincreasing.
See also \cite{Lady1968,CSTY2008}.
\end{proof}

\begin{remark}[The step that does not exist]
\label{rem:no-slide}
\(\Ga\in\Lint\) controls \(\ut=\Ga/r\) but says nothing about \(\ur\).
The strain pairing in \eqref{eq:star} contains \(\ur/r\), not \(\ut/r\).
The implication ``\(\Ga\) is bounded \(\Rightarrow\) the strain term is
controlled'' is false. No later draft may insert that step.
\end{remark}

\section{The equation for \(\Phi\)}
\label{sec:Phi-eq}

\begin{lemma}[Viscous identity for the intensive field]
\label{lem:L4}
For smooth \(\Ph\),
\begin{equation}
  \label{eq:L4id}
  \bigl(\Delta-r^{-2}\bigr)(r\Ph)=r\Lfour\Ph.
\end{equation}
\end{lemma}

\begin{proof}
\(\partial_r(r\Ph)=\Ph+r\partial_r\Ph\) and
\(\partial_{rr}(r\Ph)=2\partial_r\Ph+r\partial_{rr}\Ph\),
while \(\partial_{zz}(r\Ph)=r\partial_{zz}\Ph\). Therefore
\begin{align*}
\Delta(r\Ph)
&=2\partial_r\Ph+r\partial_{rr}\Ph
 +r^{-1}(\Ph+r\partial_r\Ph)
 +r\partial_{zz}\Ph\\
&=r\partial_{rr}\Ph+3\partial_r\Ph+r^{-1}\Ph+r\partial_{zz}\Ph.
\end{align*}
Subtract \(r^{-2}\cdot(r\Ph)=r^{-1}\Ph\):
\[
  \bigl(\Delta-r^{-2}\bigr)(r\Ph)
  =r\partial_{rr}\Ph+3\partial_r\Ph+r\partial_{zz}\Ph
  =r\Lfour\Ph.
\]
\end{proof}

\begin{lemma}[Intensive swirl equation]
\label{lem:Phi}
If \(u\) is a smooth axisymmetric solution of Navier--Stokes, then
\begin{equation}
  \label{eq:Phi}
  \partial_t\Ph+\ur\partial_r\Ph+\uz\partial_z\Ph
  +2\frac{\ur}{r}\Ph
  =\nu\Lfour\Ph.
\end{equation}
There is no extra \(-\Ph/r^2\) on the left. That term belongs to
\eqref{eq:utheta}, not to the equation for \(\Ph=\ut/r\).
On the \(\eps\)-system, the same computation produces
\begin{equation}
  \label{eq:Phieps}
  \partial_t\Ph+\ur\partial_r\Ph+\uz\partial_z\Ph
  +2\frac{\ur}{r}\Ph
  =\nu\Lfour\Ph-\eps\,\mathcal{D}^{1.3}\Ph,
\end{equation}
where \(\mathcal{D}^{1.3}\) is the restriction of \((-\Delta)^{1.3}\)
to the intensive field; the only property used later is coercivity on
\(\Hdot^{1.3}(r^3)\) via \eqref{eq:dual}.
\end{lemma}

\begin{proof}
Write \(\ut=r\Ph\) in \eqref{eq:utheta}. Then
\(\partial_t\ut=r\partial_t\Ph\),
\(\partial_r\ut=\Ph+r\partial_r\Ph\),
\(\partial_z\ut=r\partial_z\Ph\).
The left side of \eqref{eq:utheta} becomes
\begin{align*}
r\partial_t\Ph+\ur(\Ph+r\partial_r\Ph)+\uz(r\partial_z\Ph)+\ur\Ph
&=r\partial_t\Ph+r\ur\partial_r\Ph+r\uz\partial_z\Ph+2\ur\Ph\\
&=r\Bigl(\partial_t\Ph+\ur\partial_r\Ph+\uz\partial_z\Ph+2\frac{\ur}{r}\Ph\Bigr).
\end{align*}
The right side is \(\nu(\Delta-r^{-2})(r\Ph)=\nu r\Lfour\Ph\) by
Lemma~\ref{lem:L4}. Divide by \(r>0\) and extend by continuity to the
axis.
\end{proof}

\section{Pillar II: the \(\eps\)-system}
\label{sec:eps}

\begin{definition}[Fractional regularization]
\label{def:eps}
For \(\eps>0\) let \(u_\eps\) solve
\begin{equation}
  \label{eq:eps}
  \partial_t u_\eps+(u_\eps\cdot\nabla)u_\eps+\nabla p_\eps
  =\nu\Delta u_\eps-\eps(-\Delta)^{1.3}u_\eps,
  \qquad \mathrm{div}\,u_\eps=0,
\end{equation}
with smooth, compactly supported, divergence-free, axisymmetric
initial data \(u_0\). Any hyperviscous regularization strictly above
the Ladyzhenskaya--Prodi--Serrin threshold may be substituted;
the index \(1.3\) is fixed to match the existing draft.
\end{definition}

\begin{theorem}[Global smoothness for each \(\eps>0\)]
\label{thm:eps}
For every \(\eps>0\), \eqref{eq:eps} admits a unique global smooth
axisymmetric solution, and
\begin{equation}
  \label{eq:U1}
  \sup_{t\ge 0}\norm{u_\eps(t)}{L^2}^2
  +2\nu\int_0^\infty\norm{\nabla u_\eps}{L^2}^2\,\dd t
  +2\eps\int_0^\infty\norm{u_\eps}{\Hdot^{1.3}}^2\,\dd t
  \le\norm{u_0}{L^2}^2.
\end{equation}
The bound \eqref{eq:U1} does \emph{not} control
\(\int_0^T\norm{\Delta u_\eps}{L^2}^2\,\dd t\).
\end{theorem}

\begin{proof}
Take the \(L^2(\R^3)\) inner product of \eqref{eq:eps} with \(u_\eps\).
The convective term and the pressure vanish by
\(\mathrm{div}\,u_\eps=0\). The viscous term gives
\(\nu\norm{\nabla u_\eps}{L^2}^2\). The fractional term gives
\(\eps\norm{u_\eps}{\Hdot^{1.3}}^2\) by \eqref{eq:dual} with
\(\varphi=u_\eps\). This is \eqref{eq:U1}.
For each fixed \(\eps>0\) the dissipation is coercive on
\(\Hdot^{1.3}\hookrightarrow H^{1+\delta}\) in three dimensions,
which is above the Ladyzhenskaya--Prodi--Serrin line.
Galerkin approximations therefore remain smooth for all time
(Lions hyperviscous theory). Axisymmetry is preserved by uniqueness
of the cylindrical Fourier mode. Uniqueness in the smooth class is
the standard energy-difference argument.
\end{proof}

\begin{remark}[No \(H^2\) from energy]
\label{rem:no-H2}
Leray--Hopf / \eqref{eq:U1} gives
\(\nu\int_0^T\norm{\nabla u}{L^2}^2\,\dd t\le\tfrac12\norm{u_0}{L^2}^2\).
It does not give \(\int_0^T\norm{\Delta u}{L^2}^2\,\dd t\).
Any estimate that consumes \(H^2\) ``from the energy inequality'' is
circular: \(H^2\) control is what regularity would provide.
We do not do this.
\end{remark}

\begin{remark}[Weak limits do not move pointwise bounds]
\label{rem:weak}
From \eqref{eq:U1}, \(\{u_\eps\}\) is bounded in
\(L^\infty_t L^2_x\cap L^2_t H^1_x\cap L^2_t\Hdot^{1.3}_x\).
The equation gives \(\partial_t u_\eps\) bounded in
\(L^{4/3}_{\mathrm{loc}}([0,T]; H^{-s})\) for some \(s>0\)
(convective term in the usual interpolation; pressure recovered by
the Leray projector). Aubin--Lions therefore yields, up to subsequence,
strong convergence \(u_\eps\to u\) in
\(L^2_{\mathrm{loc}}([0,T]\times\R^3)\) and a Leray--Hopf limit.
Strong \(L^2\) is enough for the convective term distributionally.
It is not enough to pass an \(\Lint\) strain bound unless that bound
is already uniform in \(\eps\).
Thus \eqref{eq:B1} must be proved uniformly in \(\eps\) \emph{before}
the limit. A bound for each fixed \(\eps>0\) is free from
Theorem~\ref{thm:eps} and worthless for the classical problem.
\end{remark}

\section{The energy identity in \(r^3\,\dd r\,\dd z\)}
\label{sec:energy}

\begin{lemma}[Five-dimensional divergence]
\label{lem:div5}
If \eqref{eq:div3} holds, then
\begin{equation}
  \label{eq:div5}
  \partial_r(r^3\ur)+\partial_z(r^3\uz)=2r^2\ur,
  \qquad
  \divfive(u)=2\frac{\ur}{r}.
\end{equation}
In particular \(u\) is not divergence-free in the \(r^3\) measure.
\end{lemma}

\begin{proof}
\(\partial_r(r^3\ur)=3r^2\ur+r^3\partial_r\ur\) and
\(\partial_z(r^3\uz)=r^3\partial_z\uz\).
From \eqref{eq:div3},
\(\ur+r\partial_r\ur+r\partial_z\uz=0\), hence
\(r^3\partial_r\ur+r^3\partial_z\uz=-r^2\ur\).
Add \(3r^2\ur\).
\end{proof}

\begin{lemma}[Viscous Dirichlet form]
\label{lem:visc}
Let \(\Ph\) be smooth, with \(\Ph\) and \(\partial_r\Ph\) bounded as
\(r\to 0\), and with
\(r^3\Ph\partial_r\Ph\to 0\) as \(r\to\infty\) after \(z\)-integration,
and \(\Ph\partial_z\Ph\to 0\) as \(\abs{z}\to\infty\) after \(r\)-integration.
Then
\begin{equation}
  \label{eq:visc}
  \int(\Lfour\Ph)\,\Ph\,r^3\,\dd r\,\dd z
  =-\int\abs{\gradfive\Ph}^2\,r^3\,\dd r\,\dd z.
\end{equation}
\end{lemma}

\begin{proof}
By \eqref{eq:L4div},
\begin{align*}
\int(\Lfour\Ph)\Ph\,r^3\,\dd r\,\dd z
&=\int\partial_r(r^3\partial_r\Ph)\,\Ph\,\dd r\,\dd z
+\int(\partial_{zz}\Ph)\Ph\,r^3\,\dd r\,\dd z\\
&=\bigl[r^3(\partial_r\Ph)\Ph\bigr]_{r=0}^{\infty}
-\int\abs{\partial_r\Ph}^2 r^3\,\dd r\,\dd z
-\int\abs{\partial_z\Ph}^2 r^3\,\dd r\,\dd z.
\end{align*}
The bracket vanishes at \(r=0\) by the factor \(r^3\)
(axis regularity keeps \(\Ph\) and \(\partial_r\Ph\) bounded)
and at infinity by the decay hypothesis. There is no axis term
\(-\int\Ph(0,z)^2\,\dd z\).
\end{proof}

Cutoffs are recorded so that \eqref{eq:star} is not an ``up to the usual''
identity. Fix \(R,L\ge 1\). Let \(\chi_R\in C^\infty([0,\infty))\) satisfy
\(0\le\chi_R\le 1\), \(\chi_R\equiv 1\) on \([0,R]\),
\(\operatorname{supp}\chi_R\subset[0,2R]\),
\(\abs{\chi_R'}\le C/R\), \(\abs{\chi_R''}\le C/R^2\).
Let \(\eta_L\in C_c^\infty(\R)\) satisfy \(\eta_L\equiv 1\) on \([-L,L]\),
\(\operatorname{supp}\eta_L\subset[-2L,2L]\),
\(\abs{\eta_L'}\le C/L\). Write \(\psi_{R,L}:=\chi_R(r)\eta_L(z)\).

\begin{proposition}[Localized energy identity]
\label{prop:local}
Along a smooth solution of \eqref{eq:Phieps},
\begin{align}
\label{eq:star-loc}
\tfrac12\ddt\int\Ph^2\psi_{R,L}\,r^3
&+\nu\int\abs{\gradfive\Ph}^2\psi_{R,L}\,r^3
+\eps\int\bigl((-\Delta)^{0.65}\Ph\bigr)^2\psi_{R,L}\,r^3\nonumber\\
&=-\int\frac{\ur}{r}\Ph^2\psi_{R,L}\,r^3
+\mathcal{R}_{\mathrm{adv}}
+\mathcal{R}_{\nu}
+\mathcal{R}_{\eps},
\end{align}
where the cutoff remainders are
\begin{align*}
\mathcal{R}_{\mathrm{adv}}
&=\tfrac12\int\Ph^2\,(u\cdot\nabla\psi_{R,L})\,r^3,\\
\mathcal{R}_{\nu}
&=-\nu\int\Ph\,(\gradfive\Ph\cdot\gradfive\psi_{R,L})\,r^3,\\
\abs{\mathcal{R}_{\eps}}
&\le C\eps\norm{\Ph}{\Hdot^{1.3}(r^3)}
\norm{\Ph(1-\psi_{R,L})}{\Hdot^{1.3}(r^3)}.
\end{align*}
In particular
\[
\abs{\mathcal{R}_{\mathrm{adv}}}
+\abs{\mathcal{R}_{\nu}}
\le C\Bigl(\frac1R+\frac1L\Bigr)
\Bigl(\norm{\Ph}{L^2(A_{R,L})}^2
+\norm{u}{L^\infty(A_{R,L})}\norm{\Ph}{L^2(A_{R,L})}^2
+\norm{\gradfive\Ph}{L^2(A_{R,L})}\norm{\Ph}{L^2(A_{R,L})}\Bigr),
\]
with \(A_{R,L}\) the support of \(\nabla\psi_{R,L}\)
(the annulus \(R\le r\le 2R\) or the slabs \(L\le\abs{z}\le 2L\)).
\end{proposition}

\begin{proof}
Multiply \eqref{eq:Phieps} by \(\Ph\psi_{R,L}r^3\) and integrate.
The time term is immediate.
For advection,
\(\int(u\cdot\nabla\Ph)\Ph\psi r^3
=\tfrac12\int u\cdot\nabla(\Ph^2)\,\psi r^3\).
Integrating by parts (the cutoff has compact support, so no boundary
at infinity) and using Lemma~\ref{lem:div5},
\begin{align*}
\int u\cdot\nabla(\Ph^2)\,\psi r^3
&=-\int\Ph^2\psi\bigl(\partial_r(r^3\ur)+\partial_z(r^3\uz)\bigr)
-\int\Ph^2(u\cdot\nabla\psi)\,r^3\\
&=-2\int\frac{\ur}{r}\Ph^2\psi r^3
-\int\Ph^2(u\cdot\nabla\psi)\,r^3.
\end{align*}
Thus
\(\int(u\cdot\nabla\Ph)\Ph\psi r^3
=-\int(\ur/r)\Ph^2\psi r^3
-\tfrac12\int\Ph^2(u\cdot\nabla\psi)\,r^3\).
The stretching term contributes
\(+2\int(\ur/r)\Ph^2\psi r^3\).
The net first-order contribution on the left is
\(+\int(\ur/r)\Ph^2\psi r^3
-\tfrac12\int\Ph^2(u\cdot\nabla\psi)\,r^3\).
Move the first summand to the right; the second is
\(\mathcal{R}_{\mathrm{adv}}\).

For the viscous term,
\(\int(\Lfour\Ph)\Ph\psi r^3
=-\int\abs{\gradfive\Ph}^2\psi r^3
-\int\Ph(\gradfive\Ph\cdot\gradfive\psi)\,r^3\),
which is the Dirichlet form of Lemma~\ref{lem:visc} plus
\(\mathcal{R}_{\nu}/\nu\).
The fractional remainder is the commutator of
\((-\Delta)^{1.3}\) with multiplication by \(\psi_{R,L}\), controlled
by \eqref{eq:dual} on the support of \(1-\psi_{R,L}\).
The size of \(\mathcal{R}_{\mathrm{adv}}\) and \(\mathcal{R}_{\nu}\)
follows from \(\abs{\nabla\psi_{R,L}}\le C(R^{-1}+L^{-1})\) and Cauchy--Schwarz.
No appeal to \(\abs{f}\le\norm{f}{L^2}\) is used.
\end{proof}

\begin{theorem}[Headline energy identity]
\label{thm:star}
Let \(\Ph\) be a smooth solution of \eqref{eq:Phieps} with
\(\Ph\in L^\infty_t L^2(r^3)\cap L^2_t\Hdot^1(r^3)\) on a time interval
where the tails of Proposition~\ref{prop:local} vanish as
\(R,L\to\infty\) (in particular, for compactly supported smooth data
evolved smoothly up to time \(t\)). Then
\begin{equation}
  \label{eq:star}
  \tfrac12\ddt\norm{\Ph}{L^2(r^3)}^2
  +\nu\norm{\gradfive\Ph}{L^2(r^3)}^2
  +\eps\norm{\Ph}{\Hdot^{1.3}(r^3)}^2
  =-\int\frac{\ur}{r}\,\Ph^2\,r^3\,\dd r\,\dd z.
\end{equation}
\end{theorem}

\begin{proof}
Send \(R,L\to\infty\) in \eqref{eq:star-loc}.
The remainders vanish by the tail hypothesis and by
\(\norm{\Ph(1-\psi_{R,L})}{\Hdot^{1.3}(r^3)}\to 0\).
What remains is \eqref{eq:star}.
\end{proof}

\begin{remark}[The \(r\,\dd r\,\dd z\) comparison]
The same computation in \(r\,\dd r\,\dd z\) produces an extra axis
integral \(-\int\Ph(0,z)^2\,\dd z\) because \(\Lfour\) is then not the
Dirichlet form. The \(r^3\) calculus removes that term rather than
estimating it. No physics has left \(\R^3\).
\end{remark}

\section{The barrier}
\label{sec:B1}

\begin{proposition}[Gronwall if and only if strain]
\label{prop:gronwall}
If
\begin{equation}
  \label{eq:B1}
  \sup_{\eps\in(0,1]}
  \int_0^T\norm{\ur_\eps(t)/r}{\Lint}\,\dd t<\infty,
\end{equation}
then \eqref{eq:star} and Gronwall bound
\(\norm{\Ph_\eps(t)}{L^2(r^3)}\) uniformly in \(\eps\) on \([0,T]\).
Conversely, the only term on the right of \eqref{eq:star} is the strain
pairing, so no \(\eps\)-uniform \(\Ph\)-energy follows from
\eqref{eq:star} without a bound of type \eqref{eq:B1} or a genuinely
weaker replacement that still controls that pairing.
\end{proposition}

\begin{proof}
Let \(E_\eps(t)=\norm{\Ph_\eps(t)}{L^2(r^3)}^2\).
Then \eqref{eq:star} and H\"older,
\(\abs{\int(\ur/r)\Ph^2 r^3}\le\norm{\ur/r}{\Lint}E_\eps\),
give
\(\dot E_\eps\le 2\norm{\ur_\eps/r}{\Lint}E_\eps\),
hence
\(E_\eps(t)\le E_\eps(0)\exp\bigl(2\int_0^t\norm{\ur_\eps/r}{\Lint}\bigr)\).
Uniformity in \(\eps\) is exactly \eqref{eq:B1}.
The pointwise bound \(\abs{f(x)}\le\norm{f}{L^2}\) is false in general
and is not used: the pairing is estimated in \(\Lint\times L^1\), not
by collapsing a function to its \(L^2\) norm at a point.
\end{proof}

\begin{mdframed}[linewidth=1.0pt]
\textbf{Open problem (B1, the barrier).}
Prove \eqref{eq:B1} for arbitrary smooth axisymmetric-with-swirl data,
or replace it by a strictly weaker hypothesis that still closes
\eqref{eq:star} uniformly in \(\eps\).
This is equivalent in difficulty to global regularity of the
axisymmetric-with-swirl Navier--Stokes equations.
It is not one interpolation short of the theorem. It is the theorem.
The papers that obtain regularity in this class assume a bound of this
type \cite{CSTY2008,CSTY2009,ChenFangZhang2017}.
The assumption has not been removed.
\end{mdframed}

\subsection{Scaling: the pairing is critical}

Let \(S:=\ur/r\), \(E:=\norm{\Ph}{L^2(r^3)}^2\),
\(D:=\norm{\gradfive\Ph}{L^2(r^3)}^2\).
The pairing is \(I:=\int S\Ph^2\,r^3\).
The Navier--Stokes scaling
\(u_\lambda(x,t)=\lambda u(\lambda x,\lambda^2 t)\)
acts on components by
\(\ur_\lambda(r,z,t)=\lambda\ur(\lambda r,\lambda z,\lambda^2 t)\)
and likewise for \(\ut\). Therefore
\begin{align*}
S_\lambda(r,z,t)
&=\frac{\ur_\lambda}{r}
=\lambda^2 S(\lambda r,\lambda z,\lambda^2 t),\\
\Ph_\lambda(r,z,t)
&=\frac{\ut_\lambda}{r}
=\lambda^2\Ph(\lambda r,\lambda z,\lambda^2 t).
\end{align*}
The meridional measure transforms by
\(r^3\,\dd r\,\dd z\mapsto\lambda^{-5}\)
in the scaled variables \((\rho,\zeta)=(\lambda r,\lambda z)\).
The gradient costs an extra \(\lambda\), so
\(\abs{\gradfive\Ph_\lambda}^2\) contributes \(\lambda^6\).
Hence
\[
  I_\lambda=\lambda^{2+4-5}I=\lambda\,I,
  \qquad
  D_\lambda=\lambda^{6-5}D=\lambda\,D.
\]
The pairing and the dissipation have the same degree. The natural
Young split \(\abs{I}\le\norm{S}{\Lint}E\) returns Gronwall and
\eqref{eq:B1}. Any attempt to hide a piece in \(\nu D\) is critical
on the five-dimensional lift: an inequality
\(\abs{I}\le\delta\nu D+C_\delta(\cdots)\)
cannot send \(\delta\to 0\) independently of the size of \(S\) without
reintroducing \(\norm{S}{\Lint}\) or a supercritical norm of \(\Ph\).
The constant in front of \(D\) is of size \(1\), not size \(\delta\).
There is nothing to absorb.

A 3D Littlewood--Paley factor \(\norm{\nabla u_j}{\Lint}\lesssim 2^j\norm{u_j}{L^2}\)
is false. Bernstein in three dimensions gives \(2^{5j/2}\). The missing
\(2^{3j/2}\) has killed a previous draft and is not used here.

\begin{remark}
Any proposed closure that produces a small constant by interpolation
alone has an error. Check the scaling first.
A cubic inequality \(E'+c\nu D\le C\nu^{-3}E^3\) is supercritical.
The comparison ODE \(y'=C\nu^{-3}y^3\) explodes for large \(y(0)\).
That estimate is not a global bound and not a near-miss.
\end{remark}

\subsection{What would count as progress}

In decreasing order of difficulty:
\begin{enumerate}[leftmargin=1.2em]
\item a genuinely subcritical bound on the strain pairing for arbitrary data;
\item \eqref{eq:B1} under a hypothesis strictly weaker than those in
\cite{CSTY2008,CSTY2009,ChenFangZhang2017} (publishable on its own);
\item a sharp example showing the critical pairing cannot be absorbed.
\end{enumerate}
None of these is claimed here.

\section{Conditional continuation}

\begin{theorem}[Conditional continuation]
\label{thm:cond}
Let \(u_0\) be smooth, compactly supported, divergence-free, and
axisymmetric with swirl. Let \(u_\eps\) solve \eqref{eq:eps}.
If \eqref{eq:B1} holds on \([0,T]\), then
\(\{\Ph_\eps\}_{\eps\in(0,1]}\) is bounded in
\(L^\infty_t L^2(r^3)\cap L^2_t\dot H^1(r^3)\),
and any Aubin--Lions limit is a smooth axisymmetric solution of
classical Navier--Stokes on \([0,T]\).
If \eqref{eq:B1} holds for every finite \(T\), the classical solution
exists globally.
\end{theorem}

\begin{proof}
Proposition~\ref{prop:gronwall} gives the \(\Ph\)-bound.
Theorem~\ref{thm:gamma} and \eqref{eq:U1} remain available.
Under \eqref{eq:B1} the right-hand side of \eqref{eq:star} is uniformly
integrable, so the \(\Ph\)-energy passes to the limit.
Remark~\ref{rem:weak} gives a Leray--Hopf limit.
The swirl continuation criteria of
\cite{CSTY2008,CSTY2009}
then upgrade the limit to a smooth solution on \([0,T]\), because
\eqref{eq:B1} is precisely a bound of the type those papers assume.
\end{proof}

This is the known shape of the theory.
The content of the paper is the identity \eqref{eq:star} and the
statement that \eqref{eq:B1} is the whole remaining problem.

\section{What is not in this paper}

\begin{itemize}[leftmargin=1.2em]
\item A proof of \eqref{eq:B1}.
\item A claim that \(\Ph\)-renormalization bypasses Hardy and therefore
yields regularity. The identity bypasses one weight; the strain
returns the difficulty.
\item Gronwall-free convergence \(u_\eps\to u\) on a possibly singular
interval.
\item Any appeal to \(\int\norm{\Delta u}{L^2}^2\) from \eqref{eq:U1}.
\item Bernstein \(\norm{\nabla u_j}{\Lint}\lesssim 2^j\norm{u_j}{L^2}\)
in three dimensions (the factor is \(2^{5j/2}\)).
\item A universal constant \(6/\pi^2\), or any other constant, relating
coprime density to shell energies. There is no such theorem here.
\item An \((A,W)\) or axis-Hardy closure in the sense of \cite{ShahJun}.
\item Prime-lattice, \(Q_6\), Triple Lock, spectral-clock, or Clay
packaging.
\item Approval from a review panel of named mathematicians.
Nobody on any such list has refereed this note.
\item A declaration that the problem is closed.
\end{itemize}

\section{Related work}

Without swirl: \cite{Lady1968,UY1968}.
Partial swirl criteria: \cite{ChaeLee2002,HouLi2008,Danchin2007,ChenFangZhang2017}.
Type~I / \(\Lint\) criteria:
\cite{CSTY2008,CSTY2009,LeiZhang2011,Seregin2012}.
The May 2026 deposit is \cite{SimonsPhi2026}, with the short companion
\cite{SimonsPhiShort2026}.
Independent 2026 announcements: \cite{ShahApr,ShahMay,ShahJun}.
Those papers use, among other devices, an \((A,W)\) circulation-gradient
pair and an axis Hardy formula for \(\Ga\).
If a later independent proof of \eqref{eq:B1} is obtained on the
variables of Sections~\ref{sec:Phi-eq}--\ref{sec:energy}, it may be
stated as independent, with the May 2026 Zenodo record as antecedent
for the \(\Ph\)-identity only.

\section*{Acknowledgments}
This note corrects and restricts \cite{SimonsPhi2026}.
It does not replace that public timestamp.
It does not claim that the classical axisymmetric-with-swirl problem
is solved.

\begin{thebibliography}{99}

\bibitem{Lady1968}
O.~A.~Ladyzhenskaya,
\textit{Unique global solvability of the three-dimensional Cauchy problem
for the NS equations in the presence of axial symmetry},
Zap.\ Nauchn.\ Sem.\ LOMI \textbf{7} (1968), 155--177.

\bibitem{UY1968}
M.~R.~Ukhovskii and V.~I.~Yudovich,
\textit{Axially symmetric flows of ideal and viscous fluids filling the
whole space}, J.\ Appl.\ Math.\ Mech.\ \textbf{32} (1968), 52--62.

\bibitem{ChaeLee2002}
D.~Chae and J.~Lee,
\textit{On the regularity of the axisymmetric solutions of the NS equations},
Math.\ Z.\ \textbf{239} (2002), 645--671.

\bibitem{HouLi2008}
T.~Y.~Hou and C.~Li,
\textit{Dynamic stability of the 3D axisymmetric NS equations with swirl},
Comm.\ Pure Appl.\ Math.\ \textbf{61} (2008), 661--697.

\bibitem{Danchin2007}
R.~Danchin,
\textit{Axisymmetric incompressible flows with bounded vorticity},
Russian Math.\ Surveys \textbf{62} (2007), 73--94.

\bibitem{CSTY2008}
C.-C.~Chen, R.~M.~Strain, T.-P.~Tsai, and H.-T.~Yau,
\textit{Lower bound on the blow-up rate of the axisymmetric Navier--Stokes equations},
Int.\ Math.\ Res.\ Not.\ (2008).

\bibitem{CSTY2009}
C.-C.~Chen, R.~M.~Strain, T.-P.~Tsai, and H.-T.~Yau,
\textit{Lower bound on the blow-up rate of the axisymmetric Navier--Stokes equations~II},
Comm.\ Partial Differential Equations \textbf{34} (2009), 203--232.

\bibitem{LeiZhang2011}
Z.~Lei and Q.~S.~Zhang,
\textit{A Liouville theorem for the axially-symmetric Navier--Stokes equations},
J.\ Eur.\ Math.\ Soc.\ \textbf{13} (2011).

\bibitem{ChenFangZhang2017}
H.~Chen, D.~Fang, and T.~Zhang,
\textit{Regularity of 3D axisymmetric Navier--Stokes equations},
Discrete Contin.\ Dyn.\ Syst.\ \textbf{37} (2017), 1923--1939.

\bibitem{Seregin2012}
G.~Seregin,
\textit{A certain necessary condition of potential blow up for Navier--Stokes equations},
Comm.\ Math.\ Phys.\ \textbf{312} (2012), 833--875.

\bibitem{SimonsPhi2026}
J.~Simons,
\textit{Phi-Renormalization and Global Regularity for Axisymmetric-with-Swirl Navier--Stokes},
Zenodo, 27 May 2026.
\href{https://doi.org/10.5281/zenodo.20405404}{10.5281/zenodo.20405404};
file DOI
\href{https://doi.org/10.5281/zenodo.20405405}{10.5281/zenodo.20405405}.

\bibitem{SimonsPhiShort2026}
J.~Simons,
\textit{Phi-Renormalization and Axisymmetric-with-Swirl Navier--Stokes:
Algebraic Cancellation of the Axis Singularity},
Zenodo, 27 May 2026.
\href{https://doi.org/10.5281/zenodo.20405597}{10.5281/zenodo.20405597}.

\bibitem{ShahApr}
R.~Shahmurov,
\textit{Unconditional Axis-Regularity in the 5D Corridor},
arXiv:2604.03519.

\bibitem{ShahMay}
R.~Shahmurov,
\textit{Large-Data Global Regularity for Three-Dimensional Navier--Stokes~I},
arXiv:2605.01875.

\bibitem{ShahJun}
R.~Shahmurov,
\textit{Global Regularity for Axisymmetric Navier--Stokes Flows with Swirl},
arXiv:2606.07869.

\end{thebibliography}

\end{document}
